\section{Early Reasoning: Diferencias entre las referencias RLP, Quiet-STaR, RPT y RLPT}

Hacer que el modelo "piense antes de hablar" no fue una innovación exclusiva de RLP. Tanto Quiet-STaR como Reinforcement Pre-Training (RPT) y Reinforcement Learning on Pre-Training Data (RLPT) anticipan el razonamiento. Lo verdaderamente importante para comparar es dónde proviene la recompensa, en qué nivel de granularidad se calcula, si requiere un verificador externo y si la trazabilidad del razonamiento obtiene crédito directamente.

\subsection{De Quiet-STaR al preentrenamiento por refuerzo}

La página de las referencias coloca a Quiet-STaR, RPT y RLPT junto con RLP en una misma línea temporal. Quiet-STaR genera y evalúa el pensamiento dentro del texto ordinario; RPT convierte la corrección del siguiente token en una recompensa de aprendizaje por refuerzo verificable; RLPT utiliza datos de preentrenamiento para construir tareas de refuerzo a nivel de fragmentos más largos; mientras que RLP usa la ganancia en log-likelihood derivada del pensamiento como señal intrínseca densa para el siguiente token real.

\lecturefigure{slide-038.jpg}{Ruta de early reasoning: Quiet-STaR, RPT, RLPT y RLP}{Grabación oficial de Stanford, 00:42:20--00:43:13}

\begin{knowledgebox}{Tabla de términos: qué resuelven los métodos adyacentes}
\begin{tabularx}{\textwidth}{L{0.18\textwidth}X X}
\toprule
Método & Mecanismo central & Diferencia clave con RLP \\
\midrule
NTP & Maximización directa de la verosimilitud del siguiente token real & Sin acción de pensamiento explícita / recompensa \\
Quiet-STaR & Muestreo de pensamientos en el texto y medición de su utilidad & Enfoque temprano de auto-pensamiento, los detalles de entrenamiento difieren de RLP \\
RPT & Convertir el razonamiento del siguiente token en una tarea de RL verificable & Uso de filtrado y recompensa binaria escasa (sparse) de corrección \\
RLPT & Realizar RL a nivel de segmento sobre datos de preentrenamiento & Diferencias en la granularidad de la recompensa y el diseño del verificador \\
RLP & Recompensar el rollout con la ganancia de log-likelihood aportada por el pensamiento & Densa, a nivel de token, contrafactual sin pensar (no-think) \\
\bottomrule
\end{tabularx}
\end{knowledgebox}

\subsection{Tres ejes de comparación: fuente de recompensa, granularidad y emergencia}

NTP no tiene una recompensa separada; RPT/RLPT utilizan señales externas o verificables, generalmente más escasas; la recompensa de RLP proviene de la diferencia entre el predictor razonado interno y la línea base EMA, tomando valores continuos en la posición del token. El curso etiqueta a RLP como emergencia explícita/fuerte de razonamiento, lo cual es una generalización sobre las señales de entrenamiento por parte del autor, no una medición directa de los procesos cognitivos internos.

\lecturefigure{slide-039.jpg}{Comparación de recompensa y razonamiento entre NTP, RPT/RLPT y RLP}{Grabación oficial de Stanford, 00:43:13--00:43:56}

\begin{warningbox}{Definición precisa de Verifier-free}
RLP no requiere un modelo de determinación de respuestas adicional ni un verificador con anotaciones humanas; sin embargo, sigue utilizando el siguiente token real del corpus como supervisión y emplea una línea base EMA. Por tanto, es más preciso decir "libre de verificador externo" en lugar de "sin supervisión o sin referencia".
\end{warningbox}

\subsection{Comparación con RPT emparejada: leer junto con el mecanismo y los números}

Una tabla cualitativa solo ilustra las diferencias de diseño, pero no indica qué método es más eficaz bajo presupuestos similares; por ello, el curso sigue comparando RPT y RLP utilizando un entorno emparejado (matched setting) que sitúa la mecánica de recompensa y los resultados cuantificados en la misma página. Bajo configuraciones similares como Qwen3-1.7B-Base, OmniMath y 170 millones de tokens, el curso reporta que RLP supera a RPT con un promedio del 4\%. La explicación del mecanismo es: RPT primero utiliza un filtro externo para encontrar posiciones viables y luego otorga una recompensa binaria escasa; en cambio, RLP puede calcular la ganancia de información densa en cualquier token seleccionado y vincular el crédito a la trazabilidad del razonamiento.

\lecturefigure{slide-040.jpg}{Comparación en entorno emparejado entre RLP y RPT}{Grabación oficial de Stanford, 00:43:56--00:45:56}

\begin{knowledgebox}{Lectura de la figura: qué variables aún pueden influir en la comparación}
Aunque el modelo, los datos y el número de tokens sean similares, el prompt, el número de rollouts, la longitud del pensamiento, el optimizador, el recorte KL/clipping, el filtrado de muestras y el cómputo pueden diferir. La figura respalda que "RLP funciona mejor en los experimentos replicados por los autores", pero esto no equivale a una dominación teórica sobre todas las implementaciones de RPT.
\end{knowledgebox}

\subsection{Resumen de la sección}

RLP pertenece a la familia de early reasoning / preentrenamiento por refuerzo; su distintividad proviene del contrafactual sin pensar (no-think counterfactual) y de la recompensa basada en ganancia de información densa. Al comparar métodos adyacentes, es obligatorio examinar simultáneamente la fuente de recompensa, la granularidad, los costos de filtrado y el cómputo, en lugar de limitarse a comparar nombres de métodos o promedios individuales.
